% ==========================================
% INTRODUCTION
% ==========================================
\chapter*{Introduction}
\addcontentsline{toc}{chapter}{Introduction}

La transformation numérique des entreprises s'est imposée comme un enjeu stratégique majeur depuis le début des années 2010. Toutes les grandes organisations ont dû revoir leurs systèmes d'information, leurs infrastructures de données et leurs méthodes de travail pour faire face à l'accélération technologique et à la multiplication des sources. Dans le secteur de l'assurance, cet impératif est renforcé par la norme réglementaire, la concurrence et la nécessité de traiter des volumes de données croissants avec précision et rapidité.

C'est dans ce contexte qu'Allianz France, l'une des principales compagnies d'assurance en France, a entrepris un programme de modernisation de ses systèmes de traitement de données. Historiquement, les datamarts P\&C (Property \& Casualty) de l'entreprise étaient développés sous SAS, un logiciel propriétaire largement répandu dans les grandes institutions financières et assurantielles. Face aux limites de cette technologie (coûts de licences élevés, performances insuffisantes pour les volumes actuels et expertise de plus en plus rare) et à l'émergence de solutions open source performantes, Allianz France a lancé le projet \textbf{SAS EXIST} (SAS Exit Strategy), dont l'objectif est la migration progressive de ses datamarts vers une architecture Python/PySpark sur l'Allianz Data Platform (ADP).

C'est dans le cadre de ce projet stratégique que s'inscrit cette mission de six mois au sein du pôle Advanced Data \& Climate (ADC) de la direction P4D d'Allianz France. La mission confiée consiste à réaliser le premier PoC de cette migration, ciblant le datamart \textbf{Construction} (le marché de l'assurance construction), sélectionné pour sa représentativité et sa complexité maîtrisable en tant que cas pilote. Ce travail vise deux objectifs complémentaires : démontrer la faisabilité technique de la migration en assurant la parité fonctionnelle avec les sorties SAS de référence, et établir un cadre méthodologique et documentaire réutilisable pour les migrations à venir.

Ce rapport retrace l'ensemble du travail accompli durant cette mission. Il est organisé en cinq chapitres, le premier situe le contexte du projet SAS EXIST, le deuxième décrit les systèmes SAS existants et les technologies cibles retenues, le troisième présente la méthodologie adoptée, articulée en cinq phases, le quatrième analyse les résultats obtenus pipeline par pipeline et le cinquième propose une réflexion critique sur le projet et ouvre sur les perspectives pour la suite du programme.

\cleardoublepage
